% ======================================================================
\section{Introduction}\label{sec:intro}
% ======================================================================

\subsection{Background and motivation}

Fused Deposition Modeling (FDM) is today the most widespread additive
manufacturing process for thermoplastics. In its standard form, a
heated nozzle extrudes a continuous filament that is deposited along a
toolpath, layer by layer, until the part is complete. Virtually every
commercial FDM machine and every mainstream slicing package
(Cura, PrusaSlicer, Slic3r, \dots) shares one structural assumption: the
print head moves on a three-axis Cartesian gantry, the nozzle axis
stays vertical, and every layer is a \emph{horizontal plane} parallel to
the build plate. The slicer reduces the solid model to a stack of 2D
regions; the kinematics of the printer reduces the tool orientation to
a single fixed pose.

This assumption is a software and kinematics limitation, not a physical
one. Nothing in the extrusion process requires the nozzle to be
vertical: material can be deposited on inclined or curved surfaces,
provided that (i) the tool can be oriented so that its axis is roughly
perpendicular to the target surface, and (ii) the deposited bead is
pressed against an existing substrate. A six-degree-of-freedom
industrial robot arm satisfies requirement (i) by construction: it can
place and orient an end-effector arbitrarily within its workspace,
with the repeatability needed for extrusion processes. Removing the
planar-layer constraint then unlocks concrete benefits:

\begin{itemize}
  \item \textbf{Elimination of the staircase effect.} Sloped and curved
        surfaces are no longer approximated by stacked planar slices;
        layers can follow the actual surface curvature, so the
        as-printed geometry matches the CAD geometry without the
        characteristic stair steps of planar printing.
  \item \textbf{Conformal reinforcement.} Continuous fibers can be laid
        along the actual load paths of a structural part, on curved
        surfaces, rather than being discretized into flat slices that
        cut through the mechanically optimal trajectories.
  \item \textbf{Printing on / against existing geometry.} A
        non-planar-capable robot can print support structures that
        conform to the underside of a part, or deposit material on
        pre-existing surfaces, which is the basis of hybrid
        manufacturing.
\end{itemize}

Realizing these benefits, however, requires a slicer that fundamentally
does more than a conventional one. A planar slicer only needs to output
2D coordinates and a feed rate, because the tool orientation is
constant. A robot-based non-planar slicer must output a complete
\emph{tool pose} --- position \emph{and} orientation --- at every point
of the trajectory, must guarantee that the deposited bead stays
supported by previously deposited material or by the build plate, and
must handle path topology (connections, gaps, crossings) that planar
slicers never encounter because they can always "travel" at a fixed
height.

\subsection{Project scope and objectives}

This project develops such a slicer: a Python software package, called
\texttt{slicer}, that transforms a structural finite-element model into
full six-DoF printing trajectories, ready to be simulated and executed
on a robotic cell. The application domain is the FDM printing of
\emph{continuous-fiber composite} parts: the target objects are
load-bearing structural components (the running example is a thin
C-beam) in which stiff carbon-fiber tows are embedded in a
thermoplastic matrix, and whose fiber layout is determined by a
structural optimization (compliance minimization) performed upstream.

Concretely, the slicer must:

\begin{enumerate}
  \item \textbf{Slice a solid model into non-planar layers} ---
        each layer is an isosurface of a scalar field defined on the
        volumetric mesh, so that layers naturally follow the curved
        geometry of thin-walled structural parts.
  \item \textbf{Generate the infill of each layer, aligned with the
        fiber trajectories} deposited inside that layer. This is the
        central mechanical requirement: the matrix material must support
        the fibers, and the fiber direction controls the local
        stiffness of the printed composite.
  \item \textbf{Generate contour shells} (offset perimeter passes) at a
        nozzle-width scale, on every layer, for both the part and its
        support structure.
  \item \textbf{Automatically generate a support structure (mold)}
        underneath the overhanging regions of the part, by computing a
        "shadow volume" and slicing it into printable layers with a
        grid infill, including the handling of infill crossings so
        that no bead is deposited twice.
  \item \textbf{Assemble every layer's curves into a continuous,
        physical printing path}: order the curves, connect them when the
        connection is short and legitimate, and insert perpendicular
        approach/leave moves when it is not.
  \item \textbf{Compute a full tool pose (rotation matrix) at every
        trajectory point} --- approach axis, travel axis and their
        cross product --- so that the exported trajectory can drive a
        six-axis robot.
  \item \textbf{Export the complete multi-layer trajectory} to a
        plain-text format that an external robot simulator (RoboDK) can
        ingest directly.
\end{enumerate}

\subsection{Inputs and outputs of the pipeline}

The slicer operates within a larger tool chain built around the
\texttt{oamc} library (Optimization of Anisotropic Material with
Continuous fibers). The data flow is:

\begin{itemize}
  \item \textbf{Input: an APDL (ANSYS) finite-element model}, parsed by
        \texttt{oamc.integrations.ansys.parser.APDLParser}. The model
        contains (a) a volumetric solid mesh of the part (linear
        tetrahedra, hexahedra, or their quadratic variants), (b) a
        named surface mesh \texttt{"MOLD"} describing the surface the
        part will be printed against, and (c) boundary conditions and
        material data (an isotropic matrix material and a transversely
        isotropic fiber material).
  \item \textbf{Input: fiber paths.} For each layer, a set of
        continuous fiber trajectories (polylines in 3D), either
        computed by the structural optimization of \texttt{oamc}
        (\texttt{model.compute\_fibers}) and saved as CSV files, or
        provided directly. Each file \texttt{fiber\_<layer>\_<i>.csv}
        contains one polyline; the slicer loads them per layer.
  \item \textbf{Output: printing trajectories.} One text file for the
        part, one for the mold, each containing, line by line, every
        trajectory point with its Cartesian coordinates, the nine
        components of its tool-orientation matrix, a segment-type label
        and the layer index.
\end{itemize}

The pipeline is deliberately kept as a pure geometry/trajjectory
generator: it does not simulate the robot, check collisions, or compute
extrusion rates. Those responsibilities belong to the downstream robot
simulator. This separation keeps the slicer generic (any six-DoF cell
can consume its output) and testable (the output format is trivially
inspectable).

\subsection{Software ecosystem}

Three open-source libraries carry the computational load:

\begin{itemize}
  \item \textbf{PyVista} \cite{pyvista}, a Pythonic interface to the
        Visualization Toolkit (VTK) \cite{vtk}. Every geometric
        operation of the slicer --- isosurface extraction, scalar-field
        clipping, contouring, slicing, feature-edge extraction, Delaunay
        triangulation, convex hulls, point merging --- is a VTK filter
        exposed by PyVista. This choice is central to the project: the
        whole slicing machinery reduces to compositions of
        well-tested, C++-implemented mesh filters, called from Python.
  \item \textbf{NumPy} for all array-level mathematics, and
        \textbf{SciPy} \cite{scipy} for spatial indexing
        (\texttt{cKDTree}), convex hulls (\texttt{ConvexHull}), Delaunay
        triangulation (\texttt{Delaunay}) and rotations
        (\texttt{spatial.transform.Rotation}).
  \item \textbf{RoboDK} \cite{robodk} as the downstream consumer of the
        exported trajectories, for offline programming and simulation of
        the robot cell.
\end{itemize}

\subsection{Report organization}

The report is organized as follows.
Chapter~\ref{sec:literature} reviews the state of the art of non-planar
slicing and multi-axis printing. Chapter~\ref{sec:methodology} presents
the methodology: the complete algorithm, stage by stage, from model
orientation to trajectory export, with the underlying mathematics.
Chapter~\ref{sec:implementation} details the Python implementation:
the class architecture, the role of each module and method, the
performance engineering, and the connection to RoboDK.
Chapter~\ref{sec:results} presents results (to be completed).
